% =============================================================================
% CAPÍTULO 2: Objetivos del TFM
% =============================================================================
\chapterUnnumbered{Capítulo 2: Objetivos del TFM}

[En un párrafo explica brevemente los contenidos de este capítulo. Se
recomienda un tamaño de 1 a 3 páginas para este capítulo.]

\setcounter{section}{0}
\section{Objetivo general}

\begin{itemize}
  \item{} [El objetivo general es lo que pretendes hacer con el TFG. El
  último paso a realizar. Redáctalo en infinitivo y solo un verbo.]
\end{itemize}

\section{Objetivos específicos}

\begin{enumerate}
  \item{} [Los objetivos específicos son las partes en las que subdivides el
  trabajo para completar el objetivo general]
  \item{} [No uses más de 4--6.]
  \item{} [Siempre redactados en infinitivo.]
  \item{} [Los objetivos no son deseos para el futuro. Cuando terminas el
  proyecto se deben haber cumplido.]
  \item{} [Lo mejor es poner un objetivo para el marco teórico, otro para el
  metodológico, para el empírico, etc. Es decir, trocear el objetivo
  final en pequeñas metas volantes ;)]
\end{enumerate}
